\quellseite{293}
\part[III. The Aim of the Church]{III. The Aim of the Church in Religious Instruction}

\quellseite{295}
\chapter[§ 23. The Aim of the Church]{§ 23. The Aim of the Church in Religious Instruction}

K. Buchrucker, \emph{Grundlinien der kirchl. Katechetik}, 1889, pp. 64-115. *Fr. Zange, \emph{Evangelischer Religionsunterricht}, 1897, pp. 1-46. *E. Sachsse, \emph{Die Lehre v. d. kirchl. Erziehung}, 1897, pp. 301-328. *J. Bauer, \emph{Das Ziel der kirchl. Erziehung}, in \emph{Katechetische Zeitschrift}, 1898. *A. Eckert, \emph{Der erziehende Religionsunterricht in Schule u. Kirche}, 1899, pp. 1-29. F. M. Schiele, \emph{Gedanken ueber die Lehrbarkeit der Religion}, in \emph{Monatsschrift fuer kirchl. Praxis} 1901. H. v. Soden, \emph{Laesst sich Religion lehren?}, in \emph{Katechet. Ztschrft.} 1901. O. Baumgarten, \emph{Neue Bahnen}, 1903, pp. 48-50 and 81-84. M. v. Nathusius, \emph{Das Ziel des kirchl. Unterrichts}, 1903. A. Bonus, \emph{Vom Kulturwert der deutschen Schule}, 1904. J. Gottschick, \emph{Homiletik u. Katechetik}, 1908, pp. 110-115. J. Berndt, \emph{Methodik des Unterrichts i. d. ev. Religion}, 1909, pp. 15-35. *R. Kabisch, \emph{Wie lehren wir Religion?}, 1910, pp. 1-65. Chr. Burkstuemmer, \emph{Der Religionsunterricht i. d. Volksschule}, 1913, pp. 55-77. F. J. Kunze, \emph{Die christlich-religioese Einwirkung auf die Persoenlichkeit unter psychologischen Gesichtspunkten}, 1911. J. Steinbeck, \emph{Lehrbuch d. kirchl. Jugenderziehung}, 1914, pp. 63-77. F. H. Sneath, G. Hodges, and H. H. Tweedy, \emph{Religious Training in the School and Home}, 1917, pp. 35-70.

The definition of the aim in religious instruction is of fundamental importance from a theoretical as well as a practical aspect. Without a precise, clear definition of such aim, no reliable statements can be made concerning the educative means or material, or the correct method to be employed, and the whole process of religious instruction will bear an accidental or mechanical character. A precise definition of the aim automatically determines the material and the method to be used and places upon all our educational and training efforts the stamp of system, comprehensiveness, and purpose.

\quellseite{296}
When we inquire concerning the aim of religious instruction, we do not understand the term “instruction” in the narrow, conventional sense according to which instruction and discipline are correlated as the joint means to a higher end—the training of the young. However justified this distinction is in itself, we forbear making it in this connection. As must appear from not a few of the preceding sections, we accept so whole-heartedly the old principle that instruction must be character-building, as to identify religious instruction and religious training.

What, then, is the aim of the Church in giving religious instruction, or training? If we were speaking here of the whole vast range of the Church’s work of religious instruction, we might, on the basis of such passages as 1 Cor. 3:1-2; Hebr. 5:12, or Eph. 4:13 f. and 1 Thess. 5:23, define our aim as “Christian maturity” in the sense of Christian excellence or moral perfection. But the achievement of this aim is the task of preaching and pastoral care, and is often not accomplished before death. Besides, we dare not confound the immature Christians of whom the apostles speak, with our Christian youth. Here we deal with the task of instructing and training those children who have been baptized in infancy and are growing up in the Church, a duty which infant baptism has imperatively imposed upon us (cf. p. 4).

If we limit our task to the religious training of the young, it is clear that also here the object in view is a training for \emph{maturity}, but in another sense than that just discussed. The children growing up in the congregation, have, through Baptism, become members of the Church, and, therewith, of the body of Christ on earth. They are, however, still immature, dependent members; like minors in civil life, they do not as yet enjoy all the privileges nor do they perform all the duties, of the independent and mature members of the Church;
\quellseite{297}
they lack the prerequisites for the ones as well as the others as long as they have not been trained appropriately. Accordingly, the aim of religious instruction should, generally speaking, be the training for complete participation in the whole religious life of the mature congregation—a training not for Christian maturity in the sense of Christian excellence and relative moral perfection, but for \emph{such maturity as fits one for adequate participation in the life of the Church}.

Now the Church, or Christian congregation, according to its essence, is the throng of believers congregating around the means of grace, and standing, by means of these, in vital relation to Jesus Christ. The Church’s life is primarily a \emph{life of faith}. Accordingly, religious instruction must train for a \emph{life of personal faith} which spontaneously longs for the means of grace and through them lays hold of Christ so that the blessings of His fellowship may be enjoyed. Whatever our conception of infant baptism may be, a conscious life of faith is not one of its immediate results; this can be brought about only through appropriate training, especially since it must work itself out in all the different situations of life. Requiring, as we do, a training for a personal life of faith, we take issue with the Catholic conception of training, according to which the prime object is the drilling of the masses for the performance of churchly functions and for the external accommodation to the forms of religious life. Likewise we take issue with another conception of training, viz., the excessively intellectualistic conception which looks upon the Church as a mere school of pure doctrine and so often loses sight of the training of the heart for a personal relation to Jesus Christ or treats this as secondary in importance as compared with the inculcation of the formulae of correct doctrine. Any kind of training that does not aim at thoroughgoing influence upon the heart, upon the \emph{whole inner man}—his intellectual, emotional, volitional life, and
\quellseite{298}
does not strive to establish a \emph{personal relation} between the young and Christ, the Saviour and King, has no right of existence in the Evangelical Church.

The Christian congregation is conscious of being the communion of believers especially at common worship when it gathers around the Word and the Sacraments in order to be built up more and more into a spiritual temple; to express its life of faith through united praise and thanksgiving, and to manifest also outwardly its essential oneness. Accordingly it is the duty of the Church to train her youth for an intelligent and active participation in her \emph{worship}. If the life of faith can not exist apart from the Scripture, the same is true of worship through which the life of faith is expressed and strengthened. Hence, training for the intelligent and personal \emph{use of Scripture} must be intimately connected with the training for faith and worship. While it ought to be selfevident, in the church of the Reformation, that such training in the Scripture be given, this requirement has, so far, been met in a lamentably meager measure in Lutheran circles.

Some other factors merit consideration in this connection. As the communion of believers, the Church has received peculiar tasks from God. The Church is intended to be the institute of salvation for the world, bearing, in the footsteps of Christ, her Lord and Master, the gospel to whatever quarter is still without it; going in love and mercy after all those who have strayed or never been of her fold—a Samaritan inn for all the sick or forsaken. To carry out these \emph{tasks} is not easy; the rising generation must be trained for their fulfilment, it must be taught to know and to love them. The flock of God, furthermore, has a \emph{history} of its own, from which alone its condition in the present and also the peculiarity of the various Churches can be rightly understood. To this history the youth must be introduced
\quellseite{299}
if their own Church is to be dear and precious to them; if they are to be cognizant of what their Church has in common with others and, likewise, of the factors dividing her from others; if they are to be enabled to take a firm stand. The youth must also be made acquainted with some of the numerous stirring problems of the present; else we cannot expect them to stand their ground as mature members. Finally, the Church has a constitution. Whatever its character may be, she cannot maintain herself permanently without one. If the rising generation is eventually to attain to maturity and, with independent judgment, to take part in the planning and doing of the Church, a brief introduction to this constitution, at least in outline, cannot be omitted. Especially a Free Church, governing herself, will be compelled to insist upon acquaintance with her polity and constitution.

If, for the reasons given, the aim in the religious instruction of the young cannot be anything but the training of the rising generation for comprehensive and personal participation in the life of the mature congregation in all its ramifications, the training for a personal life of faith and for active participation in congregational worship is bound to occupy the foreground. The welfare of the rising generation depends upon the first-named factor which, in turn, is inseparable from the second; nor can there be any other foundation for the various manifestations of the life of the congregation. The personal life of faith, however, dare not receive so one-sided emphasis as to endanger the training for participation in the whole life of the congregation. In the Lutheran Church, particularly, hearty cooperation in the problems and functions of the Church at large, and even in the life of the individual congregation, has for this reason often been inhibited.

But the object of religious instruction and training must
\quellseite{300}
be defined with still greater accuracy, especially since we insist upon a \emph{personal} faith and a \emph{voluntary} participation in the life of the mature congregation; for the Christian congregation would cease to be the flock of the Lord, should its members cease to stand in such faith. Personal faith in Christ, or the reception of Jesus into the heart and life—the soul laying hold of Him as Saviour, willingly surrendering to Him as Lord and King, participating in the entire life of His flock—cannot be brought about by human instruction and training, but solely by God through His Holy Spirit. Therefore we must not ascribe to human activity what is purely a prerogative divine. It is true that the instrument whereby the educator means to influence the hearts of youth, is the Word of God, and that the Spirit dwells in this Word as the secret of its power. But is it not at the same time true that the Holy Spirit works faith in the hearers of the Word “where and when He pleases” (Augustana v)? The Word of God which the teacher of religious truth employs, is indeed efficacious; but it does not always unfold its power at the time when it is proclaimed: often it falls beneath the threshold of consciousness where it lies dormant until God’s own season shall have arrived when it shall reveal its inherent power divine. Again, it is true that the period of childhood and youth is the proper time for sowing, for the reason that it is the time when the inner development is most intensive. But who can guarantee that the power dwelling in the Word shall reveal itself in this period and that the Holy Spirit shall kindle faith in every one of the young souls entrusted to the teacher? Does not the experience of many parents and catechists and of many Christians in general, prove that the season of God, in the case of many, does not coincide with the time of youth and that very often the emotional stirrings during this period do not rise above the level of a transient sentimentality
\quellseite{301}
if they reach this level at all? Finally, it is true that the children whom the teacher of religion aims to train, are baptized children, who through the Sacrament have become children of God. But as surely as the teacher must never forget that the children entrusted to him have, objectively, become God’s own, he is at the same time aware that the subjective, God-ward development is often not merely arrested, but positively interrupted. Therefore, we shall not lay down as the aim of training the youth what has not been exclusively reserved for this period of life and what often, as experience shows, is not attained even though the instruction be of the most efficient kind. Should this be done nevertheless, discouragement or self-deception on the part of the catechist and, here and there, also on the part of the catechumen can hardly be obviated.

What the teacher of religion can accomplish is (1) to implant in the \emph{intellectual life} of the young souls entrusted to him all the material which the Holy Spirit in His own time will use for the purpose of awakening in their souls a personal, conscious faith, and of guiding them to a willing and all-sided manifestation of it in the ways and tasks of the mature congregation. (2) Secondly, wherever no unusual obstacles (chiefly the spirit in the home) bar the soul against the truth, he may succeed in so influencing the \emph{emotional life of the pupil} that “interest” is aroused, i.e., that the pupil feels himself at home in the materials of religious instruction and prefers these concepts and spheres of thought to any other, that he inwardly enters upon them, and is attracted by them. Thus a personal relation is established between the soul and the matter of instruction; the latter becomes its property—anchored in the depths of its consciousness and ready for service when the Spirit lays hold on the soul and quickens the soul. Or the Holy Spirit may let these subjects transcend the threshold of consciousness
\quellseite{302}
and release for action their inherent power divine, thus leading the soul creatively to a conscious life of faith and to voluntary participation in the life of the mature congregation. (3) Through the emotions, the teacher of religion may, finally, move the \emph{volitional life} of his pupils so that they entertain respect for the Word of God and for a sincere life of faith; that they take part with regularity in prayer, in worship, and in the tasks of the congregation. While all this is not yet the life of faith itself—opposition to all that is truly divine and spiritual may still lurk deep within the heart—, not a little has nevertheless been achieved for the subsequent life of the catechumens.

The teacher of religion is bound to keep in mind that the Church is the communion of true believers, and that the children entrusted to him are in truth members of it only when they have attained to personal faith and, as corollary thereto, to participation in all privileges and duties of the mature congregation; and he will choose all his educational and disciplinary measures with a view to the goal to which God unquestionably some time means to lead the soul. But when he is asked to define the aim of his own educational efforts, he will, in humility, answer as follows: \emph{The aim is (1) faithfully to imbed and anchor in the INTELLECT of the rising generation all the holy truths upon which the life of the mature congregation fundamentally is based, and by which alone it is constantly renewed, and without a knowledge of which one can not possibly participate in its entire life; (2) to stir the EMOTIONS to a vital interest in those truths; (3) to bend the WILL so that it may run in the paths in which the Holy Spirit, turning to account those truths, in His own season, leads to personal faith and to participation in the life of the mature congregation.} If in His grace God enables the teacher during the period of training, to arouse such a personal life of faith in this or that soul, he will
\quellseite{303}
render thanks for such gracious incentive to further efforts. But should he have no such experience, or should he, since he cannot unerringly judge, be unaware of such a result, he will by no means lose heart, but trust that the power of the Word which he has sown and the Holy Spirit who dwells in it with grace will some day become active and kindle faith in those who do not wilfully resist. Should anyone think that we have set too modest an aim, let him be reminded that an aim too high, unattainable under actual conditions, can have no other effect than to dishearten the sincere educator; that even the true aim stated above can be attained only through the blessing of God; and that, for the Christian educator, there can in fact be no higher incentive than the thought: I have the privilege of preparing the way for the Holy Spirit, of lending some help in His gracious work, of leading the souls of children entrusted to me in regard to intellect, emotions, and will, upon paths in which they shall be made by Him, soon or late, living members of the body of Christ. The incentive will be the greater the more vividly the teacher realizes that these souls have been entrusted to him in the days of youth, in the very period when the degree of their receptivity and tractability is such as to warrant the best possible prospects of permanent success.
